\section{Calculus on finite trees}\label{tree:section}

A power-mean insertion admits two expansions.  The logarithmic expansion
will compare pressures; an expansion with one canceled replica factor
will extract the perturbation identities.  Both follow from one
differentiation rule on trees.  We also bound their coefficients and the
probability of exceptional branching depths. Power-mean calculations for
hierarchical trial laws appear in M\'ezard--Parisi~\cite[Appendix~A]{MP01}
and Panchenko--Talagrand~\cite[Section~5]{PT04}; the latter uses cascade
weights introduced by Ruelle~\cite{Ruelle87}. Here we derive
the needed identities directly for finite sequences of arbitrary
probability kernels, without assuming an infinite hierarchy.

\subsection{Path kernels and insertions}

Fix an integer $L\geq1$ and numbers
\[
  0=m_0<m_1<\cdots<m_{L-1}<m_L=1.
\]
A path has depths $0,\ldots,L$.  Its depth-zero variable is the root.
For $1\leq d\leq L$, the probability kernel $K_d$ draws the next variable
given the whole prefix through depth $d-1$; $\mathscr F_d$ is the prefix
sigma-field.  The root distribution is fixed throughout all insertions.
In an $r$-level application, $L=r+1$: depths $1,\ldots,r$ carry
auxiliary variables, and depth $L$ carries a Gibbs spin configuration.
For the trial law \eqref{model:chain} with $r\ge1$, $\eta_d$ is the
depth-$d$ variable for $0\le d<r$, and $x$ is the depth-$r$ variable.

A finite tree $\mathcal T$ consists of a common root and finitely many
labeled paths, each ending at depth $L$.  Every internal vertex has at
least one child; in particular, all unary portions of paths are retained.
At each vertex, different children are drawn independently using the
appropriate kernel.  We write $E_{\mathcal T}$ for this sampling
expectation conditional on the root.  A split at depth $d$ means that
the last common vertex is at depth $d$.  Replica positions at depth $L$
are distinct vertices, but the values sampled there need not differ.
Thus independent spin replicas may be the same spin configuration.
Removing some leaves and the portions of their paths unused by the
remaining leaves preserves the sampling law of the remaining tree.

When the terminal variables are spin configurations
$\sigma^a\in\{-1,1\}^N$, the multioverlap on a subset $S$ of the
replica leaves is
\begin{equation}\label{tree:multioverlap}
 R_S=\frac1N\sum_{i=1}^N\prod_{a\in S}\sigma_i^a,
 \qquad R_\varnothing=1.
\end{equation}
An unspecified expectation $\E$ below includes the root and the
sampling of the indicated tree.

Here is how the path kernels arise from a partition function. Let
$P_d$, $1\le d<L$, be the auxiliary prior kernels, put
$X_{L-1}=\log Z$, and recursively set
\[
 X_{d-1}=m_d^{-1}\log P_d e^{m_dX_d}\qquad(1\le d<L).
\]
The corresponding \emph{tilted kernels} are
\begin{equation}\label{tree:original-tilt}
 K_d(dz\mid\mathscr F_{d-1})
 =e^{m_d(X_d(z)-X_{d-1})}P_d(dz\mid\mathscr F_{d-1}),
 \qquad 1\le d<L.
\end{equation}
They are probability kernels by the recursion. The final kernel $K_L$
is the Gibbs distribution of the spin configuration conditional on its
auxiliary prefix. When $L=1$, only this final kernel is present.

For a positive path function $A$, define
\begin{equation}\label{tree:Y-recursion}
  Y_L=A,\qquad
  Y_{d-1}=\bigl(K_dY_d^{m_d}\bigr)^{1/m_d},
  \qquad d=L,\ldots,1.
\end{equation}
We initially suppose that $A$ is bounded above and away from zero,
conditionally on the root.  Every expression below is then well
defined.  Denote the kernels after insertion by
\begin{equation}\label{tree:updated-kernel}
  K_d^A(dz\mid\mathscr F_{d-1})
  =\left(\frac{Y_d(z)}{Y_{d-1}}\right)^{m_d}
     K_d(dz\mid\mathscr F_{d-1}).
\end{equation}
The recursion shows that these are probability kernels.  Superscript
$A$ on a tree expectation means that these kernels are used.

\begin{lemma}[Insertion calculus]\label{tree:insertion}
Suppose a leaf partition function is iterated by power means, with
exponents $m_1,\ldots,m_{L-1}$, and $K_1,\ldots,K_{L-1}$ are its tilted
kernels.  Let $K_L$ be its terminal Gibbs kernel.  Multiplying each spin
weight by $A$ changes the iterated root logarithm by $\log Y_0$, where
\eqref{tree:Y-recursion} uses the kernels before insertion.  The new
kernels are \eqref{tree:updated-kernel}.

For an internal vertex $\beta$ of a prescribed tree, let $d_\beta$ be
its depth, let $n_\beta$ be its number of children, and put
\[
  \gamma_\beta=m_{d_\beta}-n_\beta m_{d_\beta+1}.
\]
The conditional joint density of its updated tree law is
\begin{equation}\label{tree:density}
  W_{\mathcal T}(A)
  :=\frac{dP_{\mathcal T}^A}{dP_{\mathcal T}}
  =\prod_{a\in\operatorname{Leaves}(\mathcal T)}A^a
       \prod_{\beta\in\operatorname{Int}(\mathcal T)}
            Y_\beta^{\gamma_\beta}.
\end{equation}
Here $Y_\beta$ is $Y_{d_\beta}$ evaluated on the prefix at $\beta$.

The recursion is monotone and homogeneous: if $a\leq A\leq b$, for
positive root-measurable $a,b$, then $a\leq Y_d\leq b$ at every depth;
and multiplying $A$ by a positive root-measurable constant multiplies
every $Y_d$ by that constant.
\end{lemma}

\begin{proof}
At the terminal spin integration, the partition-function ratio is
$K_LA=Y_{L-1}$, because $m_L=1$.  If $X_d$ and $X_d^A$ are the old and
new backward logarithms, substitution in the preceding power mean
gives
\[
  \exp\bigl(m_d(X_{d-1}^A-X_{d-1})\bigr)
  =K_d\exp\bigl(m_d(X_d^A-X_d)\bigr).
\]
Backward induction proves $X_d^A-X_d=\log Y_d$ and then gives
\eqref{tree:updated-kernel}; the same formula at $d=L$ is the usual
Gibbs density ratio.

Multiply the edge density ratios in \eqref{tree:updated-kernel} over
the tree.  An internal vertex at positive depth $d$ receives exponent
$m_d$ from its incoming edge and exponent $-n_\beta m_{d+1}$ from its
outgoing edges.  At the root there is no incoming edge, which agrees
with $m_0=0$.  Each terminal vertex receives exponent $m_L=1$.
This proves \eqref{tree:density}.  Monotonicity and homogeneity follow
at once by backward induction in \eqref{tree:Y-recursion}.
\end{proof}

An insertion may use additional variables that were absent from the
original model.  Enlarge each path kernel by their independent prior
kernels before applying the lemma.  The original backward logarithms
do not depend on these variables, so the original tilt leaves their
conditional prior kernels unchanged.  Additional root variables remain
quenched.

We will use differentiation of \eqref{tree:Y-recursion}.  For
$A_{\boldsymbol t}=1+\sum_{j=1}^k t_jD_j$, where the real path functions
$D_j$ are bounded, differentiation on any real neighborhood where
$A_{\boldsymbol t}$ is uniformly positive gives
\begin{equation}\label{tree:first-variation}
  \partial_{t_j}\log Y_d
  =E^{A_{\boldsymbol t}}\left[
       \left.\frac{D_j}{A_{\boldsymbol t}}\right|\mathscr F_d\right].
\end{equation}
The conditional expectation on the right follows one new path below
the specified prefix.  Indeed, differentiating one power mean gives
$\partial_{t_j}\log Y_{d-1}=K_d^{A_{\boldsymbol t}}
\partial_{t_j}\log Y_d$, and the terminal derivative is $D_j/A$.

\paragraph{Analytic bounds.}
The expressions are analytic on the complex domain
\[
 \sum_j|t_j|\|D_j\|_\infty<1,
\]
with bounds uniform over the kernels.
On a smaller closed domain with this sum at most $\rho<1$, $A$ has argument bounded
by $\vartheta=\arcsin\rho<\pi/2$ and modulus between $1-\rho$ and
$1+\rho$.  If $Y_d$ has argument in $[-\vartheta,\vartheta]$, then
$K_dY_d^{m_d}$ has argument in
$[-m_d\vartheta,m_d\vartheta]$, using the principal powers.  It is
nonzero, and taking its $1/m_d$ power gives the same argument bound
for $Y_{d-1}$.  More quantitatively, a lower modulus bound $a_d>0$
gives the lower bound
$a_{d-1}=a_d\cos(m_d\vartheta)^{1/m_d}>0$, whereas the upper modulus
bound $1+\rho$ is preserved.  These bounds justify integration of
analytic functions at each step and all subsequent differentiations.
They also bound \eqref{tree:density}, its canceled versions below,
and their derivatives on smaller domains, uniformly over the kernels
for fixed depth, exponents, tree size, and bounds on the inserted
functions.

\subsection{The two expansion rules}

Call a leaf \emph{protected} if its factor $A^a$ in the tree density is
canceled by a denominator $A^a$.  For a set $C$ of protected leaves and
a bounded parameter-independent function $G$, consider
$E_{\mathcal T}^A[G\prod_{a\in C}(A^a)^{-1}]$.
Differentiating
\eqref{tree:density} shows that
\[
 \partial_{t_j}\log W_{\mathcal T}(A)
   =\sum_{a\in\operatorname{Leaves}(\mathcal T)}\frac{D_j^a}{A^a}
      +\sum_{\beta\in\operatorname{Int}(\mathcal T)}
           \gamma_\beta\,\partial_{t_j}\log Y_\beta.
\]
In the derivative of
$E_{\mathcal T}^A[G\prod_{a\in C}(A^a)^{-1}]$, the leaf terms at
$a\in C$ cancel exactly with differentiation of their denominators.
Each remaining leaf term inserts $D_j^a/A^a$ and makes that leaf
protected.  By \eqref{tree:first-variation}, an internal-vertex term
is the same expectation after adjoining an independent updated path
below $\beta$, with an additional factor $D_j/A$ at its terminal
vertex.  This is precisely a fresh path at $\beta$, with coefficient
$\gamma_\beta$, and its new leaf is protected.  Conditional
independence of children and the projectivity of tree sampling justify
this representation even when $G$ uses the other children of $\beta$.

For a tree with $n$ leaves, telescoping gives the useful identity
\begin{equation}\label{tree:telescoping}
  \sum_{\beta\in\operatorname{Int}(\mathcal T)}
       \bigl(n_\beta m_{d_\beta+1}-m_{d_\beta}\bigr)=n.
\end{equation}
All summands are positive.  Define a probability law $Q_k$ on trees
with leaves labeled $1,\ldots,k$ as follows.  Start with a single path.
When the current tree has $n$ leaves, choose an internal vertex
$\beta$ with probability
\begin{equation}\label{tree:Q-rule}
  \frac{n_\beta m_{d_\beta+1}-m_{d_\beta}}{n},
\end{equation}
create a fresh child there, and continue with a completely new path
to depth $L$.  Give its terminal vertex label $n+1$.  Identity
\eqref{tree:telescoping} verifies the normalization.  The random tree
is chosen independently of all variables subsequently sampled on it.
For example, under $Q_2$ the two paths split at depth $d$ with probability
$m_{d+1}-m_d$, for $0\le d<L$. This split is uniform on the depths when
$m_d=d/L$.

\begin{lemma}[Logarithmic expansion]\label[lemma]{tree:log-expansion}
For a bounded real path function $D$ and $A=1+tD$,
\begin{equation}\label{tree:log-series}
  \log Y_0
   =\sum_{k\geq1}\frac{(-1)^{k-1}t^k}{k}
       \mathbb E_{Q_k}E_{\mathcal T}
          \prod_{a=1}^k D^a,
       \qquad |t|\|D\|_\infty<1.
\end{equation}
This identity is conditional on the root.  If
$|t|\|D\|_\infty\leq c<1$, the absolute value of the remainder after
order $K$ is at most
\begin{equation}\label{tree:log-tail}
  \sum_{k>K}\frac{c^k}{k}.
\end{equation}
\end{lemma}

\begin{proof}
The first derivative is, by
\eqref{tree:first-variation}, the expectation of $D/A$ on one updated
path.  Protect that leaf.  Every subsequent derivative can only add
a fresh path, because all current leaves are protected.  At order
$k$, all $k-1$ fresh coefficients are negative.  After normalization
by \eqref{tree:telescoping}, their absolute products are
$(k-1)!$ times the probabilities in $Q_k$.  Therefore
\[
  \left.\frac{d^k}{dt^k}\log Y_0\right|_{t=0}
    =(-1)^{k-1}(k-1)!\,
       \mathbb E_{Q_k}E_{\mathcal T}\prod_{a=1}^kD^a.
\]
Analyticity proves \eqref{tree:log-series}; the coefficients have
absolute value at most $\|D\|_\infty^k/k$, which proves the stated
convergence domain and \eqref{tree:log-tail}.
\end{proof}

Here is the companion rule with old replicas.  Fix a tree
$\mathcal T$ with an anchor leaf $a_*$.  Starting from all its old
paths, assign colors $1,\ldots,k$ in that order.  For the next color,
one may either
\begin{enumerate}[label=(\roman*),nosep]
  \item choose an unused old leaf other than $a_*$, with coefficient
  $1$; or
  \item create a fresh path from any internal vertex $\beta$ of the
  current tree, with coefficient
  $\gamma_\beta=m_{d_\beta}-n_\beta m_{d_\beta+1}$.
\end{enumerate}
Each fresh path uses a new child at its divergence vertex and new
vertices at every greater depth.  Every chosen leaf receives its color
and cannot be chosen again.  The current tree includes all old paths,
including unused ones, and every path already added.  An extension
history $E$ records these choices; write $a_j(E)$ for the leaf carrying
color $j$, $\mathcal T_E$ for its final tree, and $c_E$ for the product
of its coefficients.

\begin{lemma}[Canceled-anchor expansion]\label{tree:extension-expansion}
Let $H$ be a bounded function of the variables on the old tree, fixed
as the insertion parameters vary.  For
$A_{\boldsymbol t}=1+\sum_{j=1}^k t_jD_j$, with bounded real path
functions $D_j$, one has
\begin{equation}\label{tree:mixed-expansion}
  \left.\partial_{t_1}\cdots\partial_{t_k}
       E_{\mathcal T}^{A_{\boldsymbol t}}
        \frac{H}{A_{\boldsymbol t}^{a_*}}
  \right|_{\boldsymbol t=0}
   =\sum_E c_E\,
       E_{\mathcal T_E}
          \left[H\prod_{j=1}^kD_j^{a_j(E)}\right].
\end{equation}
The formula holds in any fixed order of the colors.  For identical
directions $D_j=D$, the coefficient of $t^k$ is the right-hand side
divided by $k!$.  The analytic bounds established above apply to its
left-hand side, uniformly over the path kernels.  The empty extension has weight $1$.
\end{lemma}

\begin{proof}
Begin with $C=\{a_*\}$ and $G=H$.
Apply the preceding rule successively, incorporating each new $D_j$
in $G$.  Fresh leaves are protected immediately, so the only
unprotected leaves are unused old leaves.  At zero all factors $A$
and $Y$ are $1$, proving the claimed rule and its coefficients.
\end{proof}

\subsection{Sums of extension coefficients}

We now control the signed normalization and the accumulated error in
the prescribed-tree identities used for the lower bound.

In an extension history, the anchor and the already colored paths
are called reference paths.  Constraints on the new color will only
specify its splitting depths with these paths.  In particular, a
constraint does not distinguish unused old paths with the same
specified ancestry.

To see the cancellation that must be controlled, start with two old
paths splitting at depth $d$, and let the first be the anchor. Require
one added color to split from the anchor exactly at $d$. It may use the
second old leaf, create a fresh path at the depth-$d$ vertex, or follow
the second old path and leave it at a depth $\ell>d$. The respective
coefficients give the total
\[
  1+(m_d-2m_{d+1})
       +\sum_{\ell=d+1}^{L-1}(m_\ell-m_{\ell+1})
  =m_d-m_{d+1}.
\]
On the uniform mesh $m_d=d/L$, with $\delta=L^{-1}$, this signed total
is $-\delta$, but the sum of absolute coefficients is $2+\delta$.
If instead the required split is delayed to $d+1<L$, the only allowed
choice is a fresh split at the unary depth-$(d+1)$ vertex on the anchor
path. Its coefficient is $-\delta$, so the absolute sum is now $\delta$.
The next two lemmas extend these two calculations: the first
computes signed totals, and the second controls absolute sums when
specified splits occur away from old branching depths.

\begin{lemma}[Reference-path totals]\label{tree:reference-total}
Suppose the next color must follow a specified reference prefix to
depth $d$ and then enter a child distinct from all $j\geq1$ children
there that contain reference paths.  Apart from these ancestry
constraints its continuation is unrestricted.  The sum of the
coefficients of all allowed choices is
\begin{equation}\label{tree:one-reference-total}
  m_d-jm_{d+1}.
\end{equation}
Consequently, if the anchor and the colors must have a prescribed
tree $S$, the total coefficient over all compatible extensions of
any old tree is
\begin{equation}\label{tree:K}
 \kappa(S)=\prod_{v\in\operatorname{Br}(S)}
           \prod_{j=1}^{a_v-1}(m_{d_v}-jm_{d_v+1}),
\end{equation}
where $a_v$ and $d_v$ are the child count and depth of each branching
vertex in the prescribed reference tree.  More generally, the same
one-step rule can be iterated starting from any already constructed
set of references.
\end{lemma}

\begin{proof}
At the specified divergence vertex, let $n$ be the child count in
the entire current tree.  The $j$ excluded children contain all
references passing through that vertex.  Each of the other $n-j$
children roots a subtree containing no reference.  Its terminal
vertices are therefore unused old leaves.  If such a subtree is
rooted at depth $d+1$, its permitted choices are an old leaf, with
coefficient $1$, or a fresh divergence at an internal vertex in the
subtree.  If it has $q$ leaves, telescoping within it gives
\[
  q+\sum_{\beta\text{ internal in the subtree}}
       (m_{d_\beta}-n_\beta m_{d_\beta+1})=m_{d+1}.
\]
The formula includes the case of a terminal child, when both sides
are $1$.  Adding these contributions to the fresh choice at the
specified vertex gives
$m_d-nm_{d+1}+(n-j)m_{d+1}=m_d-jm_{d+1}$.

For consistent prescribed splitting depths with all reference paths,
the deepest required common prefix determines the divergence vertex
and its excluded reference children.  Thus the one-step sum depends
only on the current reference tree, not on which old paths or fresh
paths realized it.  Sum backwards over successive choices.  At a
branching vertex of the final prescribed tree, the numbers of
excluded reference children when new child groups are first entered
are exactly $1,\ldots,a_v-1$.  Their factors give
\eqref{tree:K}, independently of the color order.
\end{proof}

To control errors after division by a small reference coefficient, we
also need an absolute bound.  Moving a target branching depth away from
all old branching depths forces a fresh unary split.  The next estimate
charges one small factor to each such vertex.

\begin{lemma}[Charging fresh splits]\label{tree:charging}
Assume $m_d=d/L$ and put $\delta=L^{-1}$.  Start with an old tree having
$q$ leaves and assign $k$ colors by the canceled-anchor extension
rule.  Prescribe the reference shape of the anchor and the colors,
possibly with additional restrictions on allowed extensions.  Suppose
$b$ specified branching vertices of that shape have depths which
are not branching depths anywhere in the old tree.  Then
\begin{equation}\label{tree:charging-bound}
   \sum_{E\text{ allowed}}|c_E|\leq C_{q,k}\delta^b.
\end{equation}
The constant is independent of $L$, the prescribed depths, and their
coincidences at different vertices.

In particular, for these $b$ vertices let
\[
 J=\prod_v\prod_{j=1}^{a_v-1}(m_{d_v}-jm_{d_v+1}).
\]
If $m_{d_v}\geq\eta>0$ at all of them, then $J\ne0$ and
\begin{equation}\label{tree:relative-charging}
   \frac{1}{|J|}\sum_{E\text{ allowed}}|c_E|
        \leq C_{q,k,\eta}.
\end{equation}
\end{lemma}

\begin{proof}
If the current tree has $s$ leaves, \eqref{tree:telescoping} says
that the total absolute coefficient of all fresh choices is $s$.
There are at most $s$ eligible old choices.  Hence the total absolute
coefficient at an unrestricted step is at most $2s\leq2(q+k)$.
At any fixed depth there are at most $s$ vertices.  Restricting a
step to fresh divergences at unary vertices of a prescribed depth
therefore gives total absolute coefficient at most $s\delta$, because
each such divergence has coefficient $m_d-m_{d+1}=-\delta$.

Consider an allowed history and one of its specified target
branching vertices.  That vertex cannot already be branching in the
old tree.  At the first step at which it became branching in the
current tree, a fresh path must therefore have diverged there while
it was unary.  The vertex may itself have been created earlier as
part of a new unary path; the same conclusion then applies.  Distinct
target vertices become branching at distinct steps, since a fresh
path creates only one branching vertex.  Assign each of the $b$
vertices to its creation step.  There are at most
$k!/(k-b)!$ possible assignments (and necessarily $b\leq k$).

For each fixed assignment, enlarge the family of histories by
discarding all constraints except that a designated step creates a
unary split at its assigned depth.  Bound its step sum by
$(q+k)\delta$ and every other step sum by $2(q+k)$.  Iterated summation
over choices gives at most
\[
 \frac{k!}{(k-b)!}(q+k)^b\bigl(2(q+k)\bigr)^{k-b}\delta^b.
\]
This proves \eqref{tree:charging-bound}, including repeated assigned
depths, because the designated steps remain distinct.

For the final assertion, the $j=1$ factor at each designated vertex
has absolute value $\delta$, whereas for $j\geq2$
\[
  |m_d-jm_{d+1}|=(j-1)m_d+j\delta\geq\eta.
\]
There are only a size-dependent number of the latter factors.  Thus
$|J|\geq\delta^b\eta^{\sum_v(a_v-2)}$, and the result follows.
\end{proof}

\subsection{The distribution of branching depths}

Still take $m_d=d/L$.  Suppress unary stretches only for the purpose
of describing a shape by a finite topology and its branching-depth
coordinates.  For $0<\eta<1/2$, call a shape \emph{regular} if every
branching depth divided by $L$ lies in $(\eta,1-\eta)$ and any two
distinct branching vertices have normalized depths differing by more
than $\eta$.  A singleton is regular.  This is the regularity used
when sampling a single tree; later we also use two copies whose
corresponding depth assignments are deliberately aligned.

\begin{lemma}[Regular shapes]\label[lemma]{tree:regular-shapes}
For every fixed $k$ there is $C_k<\infty$, independent of $L$ and
$\eta$, such that
\begin{equation}\label{tree:irregular-probability}
  Q_k\{\mathcal T\text{ is not regular}\}
        \leq C_k(\eta+L^{-1}).
\end{equation}
For any particular labeled topology with $a$ branching vertices and
any consistent assignment of their depths,
\begin{equation}\label{tree:assignment-probability}
  Q_k\{\text{that topology and depth assignment}\}
       \leq C_kL^{-a}.
\end{equation}
Deleting leaves from a regular tree produces a regular tree, after
discarding branching vertices that become unary.
\end{lemma}

\begin{proof}
Suppose the current tree has $n$ leaves.  A new branching vertex is
created exactly when the next path diverges at a unary vertex.
Each such choice has probability $\delta/n$ by
\eqref{tree:Q-rule}.  At any depth there are at most $n$ vertices;
hence for a set $B$ of depths the probability that the next path
creates a branching vertex at a depth in $B$ is at most $|B|\delta$.
At each of the first $k-1$ steps, the boundary depths and the depths
within normalized distance $\eta$ of any previous branching vertex
form a set of cardinality at most $C_k(\eta L+1)$.  If none of these
forbidden creation events occurs, the final tree is regular.
A union bound proves \eqref{tree:irregular-probability}.  Divergence
at an already branching vertex does not create a new coordinate and
requires no extra exclusion.

For \eqref{tree:assignment-probability}, restrict a specified final
labeled tree to leaves $1,\ldots,n$, retaining every unary stretch.
These restrictions determine its successive histories under
\eqref{tree:Q-rule}.  Each of its $a$ branching vertices is first
created by a unary divergence, whose probability is at most
$\delta$; every other transition probability is at most $1$.
The probability of the specified history is therefore at most
$\delta^a$.  If a topology description identifies equivalent
child orderings, their number is bounded in terms of $k$ alone,
which gives the stated version with $C_k$.
Finally, a retained branching vertex after deletion was already a
branching vertex of the original tree, at the same depth.  Its depth
coordinates are a subset of the old ones, so regularity is preserved.
\end{proof}

\begin{remark}\label{tree:averaging-remark}
For fixed $k$ there are finitely many labeled topologies.  Therefore
\eqref{tree:assignment-probability} converts an unnormalized depth
average $L^{-a}\sum$ for a topology with $a$ branching coordinates
into a bound for its contribution under $Q_k$.  The same is true for
errors on subshapes: summing over coordinates discarded by the
subshape costs at most the corresponding power of $L$.  Together
with \eqref{tree:irregular-probability}, this allows averaged estimates
on regular depth assignments to be used in the logarithmic expansion.
\end{remark}
